\section{Experimental Results}
\label{sec:results}

The system was evaluated in a sand-filled arena with five DW1000
anchors---four at ground level near the corners and one elevated---whose
coordinates were fixed either by the self-survey of
Section~\ref{sec:method} or by an overhead camera. The rover carried a
front and a rear tag separated by $L=\SI{0.527}{\meter}$, the rear tag
additionally equipped with a BNO085 IMU, together with an AprilTag
fiducial facing the ceiling. A ceiling-mounted Kinect~v2 tracked the
fiducial for ground truth; after intrinsic calibration and clicked-point
registration to the anchor frame the registration residual was
\SI{13.5}{\milli\meter}, and the marker lever arm was compensated per
sample. All estimation was performed off-board.

\subsection{Static accuracy and calibration}
Localization accuracy was first characterized with the rover stationary
at 19 surveyed positions across the arena ($\sim\!\SI{20}{\second}$ per
position, $\sim\!7\,600$ ranges). Table~\ref{tab:static} reports the
single-reading horizontal error under the two calibration stages.
Antenna-delay calibration alone leaves a floor-wide RMSE of
\SI{199}{\milli\meter}; adding the received-power correction
of~\eqref{eq:power} reduces it to \SI{94}{\milli\meter}
(Fig.~\ref{fig:calib}), the single largest improvement in the pipeline.

To quantify how much of the remaining error is irreducible by any static
correction, we evaluated an oracle that subtracts, from each position,
the mean residual measured there---the best achievable by any
time-invariant, position-dependent correction. Its residual RMSE is
\SI{72}{\milli\meter}, which therefore lower-bounds any fixed calibration
and shows that the leftover error is not a static spatial map that a
denser survey could remove. Temporal filtering nonetheless attains
\SI{60}{}--\SI{65}{\milli\meter} at rest, below this bound, because it
averages the zero-mean high-frequency component of the residual over
successive readings, leaving only a slowly varying multipath term that
the static oracle cannot capture.

\begin{table}[t]
  \centering
  \caption{Static single-reading horizontal error over 19 positions.}
  \label{tab:static}
  \begin{tabular}{@{}lccc@{}}
    \toprule
     & \textbf{RMSE} & \textbf{Median} & \textbf{95th pct.}\\
    \midrule
    Antenna-delay calibration only          & 199 & 132 & 393\\
    \;+ received-power correction            & \textbf{94} & \textbf{69} & 160\\
    Oracle (per-position constant)           & 72  & 31  & ---\\
    \bottomrule
  \end{tabular}\\[2pt]
  {\footnotesize All values in \si{\milli\meter}.}
\end{table}

\begin{figure}[t]
  \centering
  \begin{tikzpicture}
    \begin{axis}[
      width=0.98\columnwidth, height=4.5cm,
      ybar, bar width=17pt,
      ymin=0, ymax=225, xmin=0.45, xmax=2.55,
      ylabel={Single-reading 2-D RMSE (\si{\milli\meter})},
      ylabel style={font=\small},
      xtick={1,2},
      xticklabels={{Antenna\\delay only}, {+\,Power-bias\\correction}},
      xtick style={draw=none},
      x tick label style={font=\small, align=center},
      ytick={0,50,...,200}, tick label style={font=\footnotesize},
      ymajorgrids, grid style={cGrayL}, axis lines=left,
      nodes near coords, nodes near coords style={font=\small\bfseries},
      every node near coord/.append style={/pgf/number format/precision=0,
        /pgf/number format/fixed},
    ]
      \addplot[draw=cAmber!70!black, fill=cAmberL] coordinates {(1,199) (2,94)};
      \draw[dashed, cRed, line width=0.9pt] (axis cs:0.55,72) -- (axis cs:2.45,72);
      \node[font=\footnotesize, text=cRed, anchor=south] at (axis cs:1.5,74)
        {oracle floor \SI{72}{\milli\meter}};
    \end{axis}
  \end{tikzpicture}
  \caption{Effect of the two-step calibration on floor-wide single-reading
  error. The received-power correction removes the strength-dependent
  residual left by antenna-delay calibration; the dashed line marks the
  \SI{72}{\milli\meter} oracle floor of any fixed correction.}
  \label{fig:calib}
\end{figure}

\subsection{Estimator comparison and anchor self-survey}
Table~\ref{tab:est} compares the moving-horizon estimator with an
extended Kalman filter on identical recorded ranges. While the vehicle is
moving, the MHE reduces trajectory jerk by \SI{40}{}--\SI{44}{\percent}
(from 0.61 to \SI{0.36}{\meter\per\second\cubed}) for only
\SI{6}{\milli\meter} of additional lag, and the two estimators are
equivalent at rest; each MHE solve completes in \SI{3}{\milli\second},
within the \SI{6}{\hertz} update budget. The reduced jerk reflects the
window optimization rejecting isolated range excursions that a
single-step filter propagates.

The anchor self-survey was evaluated against the independently registered
overhead-camera layout (Table~\ref{tab:survey}). The raw inter-anchor
ranges are biased long by up to \SI{202}{\milli\meter}
(RMS \SI{118}{\milli\meter}) because each measurement carries the antenna
delays of both participating anchors. After multidimensional scaling and
alignment without scaling, the reconstructed constellation agrees with
the camera layout to \SI{51}{\milli\meter} RMS once the per-anchor delay
bias is removed---an independent confirmation of the geometry to within
centimetres. The deployed configuration is built from this self-survey.

\begin{table}[t]
  \centering
  \caption{Moving-horizon estimator versus Kalman filter (same logs).}
  \label{tab:est}
  \begin{tabular}{@{}lcc@{}}
    \toprule
    \textbf{Metric} & \textbf{Kalman (EKF)} & \textbf{MHE}\\
    \midrule
    Trajectory jerk, moving (\si{\meter\per\second\cubed}) & 0.61 & \textbf{0.36}\\
    Additional lag (\si{\milli\meter})       & ---  & 6\\
    RMSE at rest (\si{\milli\meter})         & 64   & \textbf{61}\\
    Solve time (\si{\milli\second})          & $<1$ & 3\\
    \bottomrule
  \end{tabular}
\end{table}

\begin{table}[t]
  \centering
  \caption{Anchor self-survey versus overhead-camera layout.}
  \label{tab:survey}
  \begin{tabular}{@{}lc@{}}
    \toprule
    \textbf{Quantity} & \textbf{Value (\si{\milli\meter})}\\
    \midrule
    Raw inter-anchor range bias (max)          & 202\\
    Raw inter-anchor range bias (RMS)          & 118\\
    Reconstructed geometry vs.\ camera (RMS)   & \textbf{51}\\
    \bottomrule
  \end{tabular}
\end{table}

\subsection{Trajectory tracking}
The full system was finally evaluated on a driven run under continuous
optical ground truth. Over the trajectory the Kinect provided 133 valid
pose samples at near-continuous coverage, and both
tags were serviced evenly (355 and 362 sweeps) with an average of five
anchors per sweep, so the ground truth constrains the entire path rather
than isolated segments. The rigid-body MHE track follows the reference
with a horizontal RMSE of \SI{124}{\milli\meter}; the estimated path
length is $1.4\times$ that of the truth, indicating that the residual is
dominated by high-frequency wander rather than a systematic offset, as
the two trajectories share the same overall shape. This moving accuracy
is consistent with the single-reading static scatter accumulating under
motion, and is limited by per-sweep range noise, which longer averaging
windows trade against lag; reducing it further through tighter range
gating and higher sweep rates is the subject of ongoing work.
